\documentclass[11pt,a4paper]{article}
\usepackage[a4paper,margin=2.5cm]{geometry}
\setlength{\parindent}{0pt}
\setlength{\parskip}{0.6em}
\usepackage{helvet}
\renewcommand{\familydefault}{\sfdefault}
\usepackage[utf8]{inputenc}
\usepackage[T1]{fontenc}
\usepackage{setspace}
\onehalfspacing
\usepackage{titlesec}
\titleformat{\section}{\bfseries\Large\centering}{}{0pt}{}
\titlespacing*{\section}{0pt}{2em}{1em}

\begin{document}

\begin{center}
\textbf{\Large REGIMENTO INTERNO}\\
\textbf{\large MOVIMENTO EM DEFESA DA EDUCAÇÃO -- MoveEduca}
\end{center}

\section*{CAPÍTULO I -- DAS DISPOSIÇÕES INICIAIS}

\textbf{Art. 1º} Este Regimento Interno disciplina o funcionamento administrativo, associativo, operacional e de governança do Movimento em Defesa da Educação, denominado abreviadamente MoveEduca, em complemento ao seu Estatuto Social.

\textbf{Art. 2º} O Regimento Interno será interpretado de acordo com o Estatuto Social, a legislação aplicável, as políticas institucionais aprovadas e a missão de despertar a consciência do ser por meio da educação.

\textbf{Art. 3º} As atividades do MoveEduca observarão seus valores institucionais de bem do todo, união, neutralidade política, respeito, colaboração, integridade, sustentabilidade, inovação e transparência.

\section*{CAPÍTULO II -- DOS ASSOCIADOS}

\textbf{Art. 4º} A admissão de associado dependerá de requerimento escrito, apresentação de dados cadastrais, ciência do Estatuto Social e deste Regimento Interno, e aprovação pela Diretoria Executiva.

\textbf{Art. 5º} O cadastro associativo conterá, no mínimo, nome completo ou razão social, nacionalidade, estado civil quando aplicável, profissão ou atividade, RG, CPF ou CNPJ, endereço, e-mail, telefone, categoria associativa, data de admissão e situação cadastral.

\textbf{Art. 6º} O associado deverá manter seus dados atualizados e comunicar alterações relevantes à Secretaria-Geral.

\textbf{Art. 7º} O associado poderá requerer desligamento voluntário por escrito, cabendo à Secretaria-Geral registrar o pedido e à Diretoria Executiva tomar ciência na reunião subsequente.

\textbf{Art. 8º} O procedimento de exclusão de associado por justa causa observará o Estatuto Social e será conduzido por escrito, com indicação dos fatos, documentos disponíveis, prazo de defesa, decisão fundamentada e possibilidade de recurso à Assembleia Geral.

\textbf{Art. 9º} Medidas disciplinares poderão consistir em advertência, suspensão temporária de direitos associativos ou exclusão, conforme a gravidade da conduta, a reincidência, o dano causado e a proporcionalidade da medida.

\section*{CAPÍTULO III -- DAS ASSEMBLEIAS E REUNIÕES}

\textbf{Art. 10.} As Assembleias Gerais serão convocadas na forma do Estatuto Social, com indicação clara da ordem do dia, data, horário, local ou meio eletrônico de participação.

\textbf{Art. 11.} Assembleias, reuniões, votações, consultas, registros de presença e deliberações poderão ocorrer por meio presencial, híbrido ou eletrônico, inclusive por aplicativo, plataforma digital ou sistema próprio.

\textbf{§ 1º} A participação eletrônica será admitida quando houver identificação dos participantes, segurança do voto, direito de participação, direito de manifestação e geração de registros auditáveis.

\textbf{§ 2º} Os registros eletrônicos poderão incluir lista de presença, relatório de sistema, extrato de votação, gravação, registro de mensagens, confirmação eletrônica ou outro meio idôneo definido pela Diretoria Executiva.

\textbf{Art. 12.} As atas deverão registrar, no mínimo, data, horário, local ou meio utilizado, forma de convocação, participantes, pauta, deliberações, votos quando necessário e assinaturas do Presidente e do Secretário da mesa.

\textbf{Art. 13.} Reuniões da Diretoria Executiva, do Conselho Fiscal e do Conselho de Administração poderão ocorrer presencialmente, remotamente ou em formato híbrido, com registro em ata, memória de reunião ou documento equivalente.

\section*{CAPÍTULO IV -- DA DIRETORIA EXECUTIVA}

\textbf{Art. 14.} A Diretoria Executiva conduzirá a gestão ordinária do MoveEduca, executando programas, projetos, contratos, parcerias, orçamento, controles internos e deliberações dos órgãos estatutários.

\textbf{Art. 15.} A Diretoria Executiva poderá distribuir atribuições operacionais entre seus membros, colaboradores, voluntários ou prestadores de serviços, sem transferir responsabilidades estatutárias.

\textbf{Art. 16.} Contratos, instrumentos de parceria, compromissos financeiros e movimentações bancárias observarão o Estatuto Social, a Política Financeira e as deliberações competentes.

\textbf{Art. 17.} O Presidente representará institucionalmente o MoveEduca e poderá delegar atos específicos por instrumento escrito, com prazo e poderes definidos.

\textbf{Art. 18.} O Diretor Financeiro manterá controles financeiros organizados, documentos comprobatórios e relatórios periódicos para a Diretoria Executiva, o Conselho Fiscal e demais órgãos competentes.

\textbf{Art. 19.} O Secretário-Geral manterá arquivo de atas, cadastros, comunicações, documentos institucionais e registros necessários ao funcionamento associativo.

\section*{CAPÍTULO V -- DOS CONSELHOS}

\textbf{Art. 20.} O Conselho de Administração atuará de forma estratégica e não executiva, orientando planejamento, governança, riscos, políticas institucionais e parcerias relevantes.

\textbf{Art. 21.} O Conselho Fiscal exercerá fiscalização econômico-financeira, contábil e patrimonial, podendo solicitar documentos, esclarecimentos e relatórios à Diretoria Executiva.

\textbf{Art. 22.} Os conselhos deverão registrar suas reuniões e recomendações, preservando independência, confidencialidade e declaração de eventuais conflitos de interesse.

\section*{CAPÍTULO VI -- DOS PROJETOS, PARCERIAS E VOLUNTARIADO}

\textbf{Art. 23.} Projetos e parcerias deverão ser compatíveis com as finalidades educacionais, sociais, culturais, tecnológicas, inclusivas e de fortalecimento institucional do MoveEduca.

\textbf{Art. 24.} Todo projeto deverá ter responsável designado, objetivos, público, cronograma, orçamento estimado quando houver recursos financeiros, forma de acompanhamento e documentos de suporte.

\textbf{Art. 25.} Parcerias com a Administração Pública observarão a Lei nº 13.019/2014, o edital, o plano de trabalho, o instrumento firmado, as normas do órgão parceiro e a prestação de contas exigida.

\textbf{Art. 26.} A atuação de voluntários deverá ser formalizada por termo próprio quando exigido ou recomendável, com definição de atividades, período, responsabilidades e ausência de vínculo empregatício.

\section*{CAPÍTULO VII -- DA CONDUTA, COMUNICAÇÃO E INTEGRIDADE}

\textbf{Art. 27.} Associados, dirigentes, conselheiros, empregados, voluntários e prestadores de serviços deverão agir com respeito, boa-fé, integridade, colaboração e compromisso com as finalidades institucionais.

\textbf{Art. 28.} É vedado usar o nome, a marca, os canais, os documentos, os recursos ou a reputação do MoveEduca para fins pessoais, partidários, discriminatórios, ilícitos ou incompatíveis com o Estatuto Social.

\textbf{Art. 29.} A comunicação institucional deverá ser clara, responsável, respeitosa, coerente com a missão e os valores da associação, e autorizada pela Diretoria Executiva ou por pessoa designada.

\textbf{Art. 30.} Situações de conflito de interesses deverão ser declaradas por escrito, com abstenção de voto ou influência na deliberação correspondente.

\section*{CAPÍTULO VIII -- DOS DOCUMENTOS E DA TRANSPARÊNCIA}

\textbf{Art. 31.} O MoveEduca manterá arquivo físico ou digital de seus documentos estatutários, atas, cadastros, contratos, termos, relatórios, prestações de contas, documentos financeiros e registros de projetos.

\textbf{Art. 32.} O acesso a documentos internos poderá ser solicitado por associado em pleno gozo de seus direitos, observadas a proteção de dados pessoais, o sigilo legal, a segurança institucional e a razoabilidade do pedido.

\textbf{Art. 33.} Relatórios anuais, demonstrações financeiras e prestações de contas serão submetidos aos órgãos competentes conforme o Estatuto Social.

\section*{CAPÍTULO IX -- DAS DISPOSIÇÕES FINAIS}

\textbf{Art. 34.} Este Regimento Interno poderá ser alterado por deliberação da Assembleia Geral, quando submetido pela Diretoria Executiva, pelo Conselho de Administração ou por associados legitimados.

\textbf{Art. 35.} Os casos omissos serão resolvidos pela Diretoria Executiva, observados o Estatuto Social, as políticas institucionais, a legislação aplicável e posterior ratificação pela Assembleia Geral quando envolver matéria de sua competência.

\vspace{2cm}
\begin{center}
\textbf{[CIDADE/UF]}, \textbf{[DATA]}.

\vspace{2cm}
\rule{8cm}{0.4pt}\\
\textbf{[NOME COMPLETO]}\\
Presidente do MoveEduca

\vspace{1.5cm}
\rule{8cm}{0.4pt}\\
\textbf{[NOME COMPLETO]}\\
Secretário(a)-Geral do MoveEduca

\vspace{1.5cm}
\rule{8cm}{0.4pt}\\
\textbf{[NOME DO ADVOGADO]}\\
\textbf{[OAB/UF Nº]}
\end{center}

\end{document}
